

En
composant (2.1.2) et (2.1.3), on obtient les homomorphismes (2.1.1).

La structure d'algèbre définie par les homomorphismes (2.1.1) est associative (cf. [1] XI 2). Par ailleurs, le @@M0@@-module @@M1@@ est muni de façon standard d'une structure de @@M2@@-module en faisant opérer @@M3@@ de façon évidente, et @@M4@@ par fonctorialité; de même, les @@M5@@ sont munis de deux structures de @@M6@@-modules. Il est clair que les homomorphismes (2.1.2) et (2.1.3) sont compatibles avec ces deux structures: pour (2.1.2), par fonctorialité (en se plaçant dans la catégorie des @@M7@@-modules), et pour (2.1.3), en utilisant le résultat classique sur les morphismes d'un complexe acyclique. Nous avons donc muni @@M8@@ d'une structure de @@M9@@-algèbre associative.

\subsection*{2.2.}
Supposons que @@M10@@ admette une résolution par des @@M11@@-modules projectifs @@M12@@, telle que @@M13@@ soit une @@M14@@-algèbre graduée, la différentielle satisfaisant:
@@M15@@
\SGAAuthorFootText{1}{On peut montrer qu'une telle résolution existe toujours.}
Alors les applications ``produit'' définissent des diagrammes commutatifs pour @@M16@@:
@@M17@@
et par suite un homomorphisme de complexes @@M18@@ qui prolonge l'homomorphisme d'augmentation : @@M19@@. Ceci montre que la structure d'algèbre de @@M20@@ peut se calculer à partir de celle de @@M21@@. En particulier, si @@M22@@ est une @@M23@@-algèbre alternée (i.e. telle que pour tout
